\documentclass[11pt,a4paper]{article}

\usepackage[utf8]{inputenc}
\usepackage[T1]{fontenc}
\usepackage{amsmath,amssymb}
\usepackage{geometry}
\usepackage{graphicx}
\usepackage{hyperref}
\usepackage{booktabs}
\usepackage{enumitem}
\usepackage{xcolor}

\geometry{margin=2.7cm}
\hypersetup{
    colorlinks=true,
    linkcolor=blue!50!black,
    citecolor=blue!50!black,
    urlcolor=blue!50!black
}

\title{\vspace{-1.5em}Justification for the Spectral Loss Design:\\
Wasserstein, KL-Divergence, and Soft-Max Peak Height}
\author{Wave Spectrum Forecasting with Transformers}
\date{\today}

\begin{document}
\maketitle
\vspace{-2em}

\section{Context and Purpose}

The model under development forecasts the directional wave spectrum
$E(f,t)$ (or its unit-area normalized form, $E(f,t)/m_0$, the \texttt{shape}
target) over a log-spaced frequency grid of 47 bins, between 0.02 and
0.485~Hz. Training and evaluation have so far relied on RMSE computed
bin-by-bin between the predicted and observed spectrum. This document lays
out the technical justification for complementing that RMSE term with three
additional, physically motivated penalty terms — a Wasserstein-1 distance, a
Kullback-Leibler (KL) divergence, and a differentiable soft-max peak-height
term — all three already implemented in this project as
\texttt{SpectralWassersteinLoss}, \texttt{SpectralKLDivergenceLoss}, and
\texttt{SoftPeakHeightLoss} (\texttt{utils/loss.py}) — and situates that
choice relative to how these mechanisms are already used in analogous
domains of spectral/distributional comparison and differentiable extremum
localization (audio, EEG, image retrieval, non-negative matrix
factorization, human pose estimation).

\section{Why RMSE Is a Poor Metric for Spectral Evaluation}
\label{sec:rmse-problem}

RMSE between two spectra is defined bin by bin:
\begin{equation}
\text{RMSE}(E_{\text{pred}}, E_{\text{true}}) =
\sqrt{\frac{1}{N}\sum_{i=1}^{N}\big(E_{\text{pred}}(f_i) - E_{\text{true}}(f_i)\big)^2}.
\end{equation}
This formulation treats each frequency bin as an independent, interchangeable
coordinate: the error between $f_i$ and $f_{i+1}$ is computed exactly the
same way as the error between $f_i$ and $f_{40}$. The frequency axis — which
has real, non-uniform, log-spaced geometry — is invisible to the metric.
This produces three concrete problems for a wave spectrum:

\begin{enumerate}[leftmargin=1.4em]
\item \textbf{No notion of proximity between bins.} A peak predicted one bin
off (a small, physically minor phase/position error) is penalized comparably
to a peak that vanishes entirely or appears in a completely wrong frequency
band (a physically severe energy/modality error). RMSE cannot distinguish
these because there is no cost for \emph{moving} energy mass from one bin to
its neighbor — only a cost for it being in the right place or not.

\item \textbf{Structural bias toward smoothed spectra.} When there is
genuine uncertainty about the exact position of the spectral peak across
training samples (normal in ocean waves — the peak frequency $f_p$ drifts
hour to hour), the estimator that minimizes mean squared error is not "predict
the most likely peak" but the \emph{average} of several narrow spectra
slightly shifted relative to one another. That average is broader and lower
than any individual sample. This is the same mechanism, well documented,
behind blur in image-generation models trained with pixelwise $L_2$/MSE loss.

\item \textbf{Insensitivity to modality errors.} In mixed sea states
(wind sea + swell), the spectrum is bimodal or trimodal. A whole-spectrum
RMSE dilutes exactly the error type that matters physically: getting total
energy right while misplacing the relative amplitude or position of the two
peaks, or collapsing two separate peaks into one.
\end{enumerate}

In short: RMSE is well suited to \emph{magnitude} accuracy per bin, but blind
to \emph{shape} and \emph{position} — precisely what matters most in a wave
spectrum, since total energy (proportional to $H_s^2$) is already tracked
separately by the bulk-parameter metrics (\texttt{Hs\_RMSE},
\texttt{Tm02\_RMSE}) already computed in this pipeline.

\section{The Wasserstein Distance as a Correction}
\label{sec:wasserstein}

The Wasserstein distance treats each spectrum (normalized to unit area,
exactly what \texttt{compute\_shape} produces) as a probability distribution
over frequency, rather than a vector of independent coordinates. For two
distributions $\mu,\nu$ over the same metric space:
\begin{equation}
W_p(\mu,\nu) = \left(\inf_{\gamma \in \Pi(\mu,\nu)}
\int |f-g|^p \, d\gamma(f,g)\right)^{1/p},
\end{equation}
where $\Pi(\mu,\nu)$ ranges over all transport plans moving the mass of
$\mu$ onto $\nu$. Intuitively, $W_p$ is the minimal cost of physically
moving energy from one part of the spectrum to another, where moving one
unit of mass a given distance costs proportionally to that distance.

For one-dimensional distributions — exactly the case here, $E(f)/m_0$ over
the frequency grid — $W_1$ has a closed form via cumulative distribution
functions (CDFs):
\begin{equation}
W_1(\mu,\nu) = \int_{0}^{\infty} \big| F_\mu(f) - F_\nu(f) \big| \, df,
\end{equation}
which is exactly what \texttt{SpectralWassersteinLoss} computes
(via a discrete cumulative-trapezoid integral over the true grid).

This directly resolves the three problems from Section~\ref{sec:rmse-problem}:
a shifted peak costs little (proportional to how far it moved), an
optimum under $W_1$ does not collapse to a blurred average the way an
$L_2$-minimizing optimum does (translated copies of the same sharp bump
have a non-blurred Wasserstein barycenter), and a modality error (mass moved
from one peak to a distant one) is directly and heavily penalized because
the transported distance is large.

\section{Why Combine RMSE with a Wasserstein Penalty}
\label{sec:hybrid}

The implemented approach does not replace RMSE — it composes:
\begin{equation}
\mathcal{L} = \text{RMSE}\big(E_{\text{pred}}, E_{\text{true}}\big)
\;+\; \lambda_W \cdot W_1\big(\hat P_{\text{pred}}, \hat P_{\text{true}}\big),
\qquad \hat P = E/m_0.
\end{equation}

Reasons to keep RMSE as the primary term rather than training on $W_1$ alone:

\begin{enumerate}[leftmargin=1.4em]
\item \textbf{Per-bin magnitude fidelity still matters.} $H_s = 4\sqrt{m_0}$
and $T_{m02}$ are computed by direct integration of the spectrum; those
integrals depend on point-wise magnitude, not only on the relative shape or
position of energy mass. RMSE anchors that component.
\item \textbf{$W_1$ is defined on normalized distributions.} By
construction it measures \emph{shape} error, not total energy — comparing
two unit-mass distributions discards scale information. This is desirable
as a complementary term, but it cannot be the sole loss for a target that
also needs absolute-magnitude accuracy.
\item \textbf{Training stability.} RMSE is smooth, strictly convex around
the target, and well-behaved from the start of training, when predictions
are still far from the target distribution.
\item \textbf{$\lambda_W$ exposes the trade-off explicitly}, tunable via
validation, between bin-wise fidelity and shape/position fidelity.
\end{enumerate}

\section{A Remaining Gap: Physically-Weighted Local Errors}
\label{sec:gap}

RMSE and $W_1$ together still share one blind spot. Both treat every
frequency bin as equally important \emph{in isolation}: RMSE (computed in
log-space, on $\log E(f)$ or $\log(E(f)/m_0)$, floored near zero to avoid
$\log(0)$) assigns the same weight to a bin crossing that near-zero floor
regardless of whether that bin is the spectral peak or an inconsequential
tail bin; $W_1$'s transport cost depends only on how far mass moved, not on
how physically important the bin it moved from or to actually is.

This matters because the bulk parameters that ultimately matter —
$H_s$, $T_{m02}$ — come from integrals ($m_0=\int E\,df$,
$m_2=\int f^2E\,df$) that are dominated by the region around the spectral
peak. A log-space error at a near-zero tail bin contributes almost nothing
to these integrals; the identical log-space error at the peak can
substantially distort them. Neither RMSE nor $W_1$, as currently combined,
makes this distinction explicit — a further term is needed that weights
each bin's error in proportion to how much of the spectrum's actual energy
that bin represents.

\section{KL Divergence as a Physically-Weighted Penalty}
\label{sec:kl}

\subsection{Definition and Derivation}

Since $\hat P = E/m_0$ is already, by construction, a non-negative,
unit-area probability distribution over frequency, the standard
information-theoretic measure of divergence between two distributions
applies directly — the Kullback-Leibler divergence:
\begin{equation}
D_{KL}(P\,\|\,Q) = \sum_i P(f_i)\,\log\frac{P(f_i)}{Q(f_i)}.
\end{equation}
Expanding the logarithm, $\log\frac{P(f_i)}{Q(f_i)} = \log P(f_i) - \log
Q(f_i)$ (no extraneous sign), and writing $\epsilon_i = \log P(f_i) - \log
Q(f_i)$ for the per-bin \emph{log-space error} — the same log-space
representation the model already predicts and trains in:
\begin{equation}
D_{KL}(P\,\|\,Q) = \sum_i P(f_i)\cdot \epsilon_i.
\end{equation}
Since $\sum_i P(f_i) = 1$, this is exactly a \textbf{$P$-weighted average}
of the per-bin log-space errors — the weight on bin $i$ is $P(i)$, the
fraction of the true spectrum's total energy actually located there.
Contrast with the plain log-space MSE currently used as the main per-bin
loss:
\begin{equation}
\mathcal{L}_{\text{MSE-log}} = \frac{1}{N}\sum_i \epsilon_i^2,
\end{equation}
where every bin carries the same weight, $1/N$, independent of position.

\subsection{Non-Negativity and Relationship to Cross-Entropy}

Individual terms $P(f_i)\cdot\epsilon_i$ can carry either sign — positive
where $Q$ underestimates $P$, negative where $Q$ overestimates it — but the
weighted sum is always non-negative, with equality if and only if $P=Q$
everywhere (Gibbs' inequality, a direct consequence of the concavity of
$\log$). This guarantees $D_{KL}$ behaves as a proper training penalty,
bounded below by zero. It is worth stating explicitly because expanding the
logarithm the wrong way — introducing a spurious overall minus sign — yields
$-D_{KL}(P\|Q)$, which is non-positive and would, if minimized, reward the
optimizer for making the prediction worse rather than better; this is a
one-symbol error easy enough to make that it is worth flagging as a check
on the final implementation, not just the derivation.

$D_{KL}$ is closely related to the cross-entropy $H(P,Q) =
-\sum_i P(f_i)\log Q(f_i)$:
\begin{equation}
D_{KL}(P\,\|\,Q) = H(P,Q) - H(P),
\end{equation}
where $H(P)$ is $P$'s own Shannon entropy. Since $H(P)$ does not depend on
the prediction $Q$, $D_{KL}$ and $H(P,Q)$ produce \emph{identical gradients}
with respect to the predicted spectrum — this is exactly the log-softmax
cross-entropy mechanism used in classification, applied over the frequency
axis instead of over class labels, and it means the choice between the two
is not a training-behavior decision. $D_{KL}$ is preferred here purely for
interpretability of the reported value: it is exactly zero at a perfect
prediction, whereas plain cross-entropy floors at $H(P)$, a nonzero,
sample-dependent value that makes "how close to perfect" harder to read
off the raw loss.

\subsection{Physical Interpretation for Bulk Wave Parameters}

For the \texttt{shape} target, $P(f_i) = \big(E(f_i)/m_0\big)\cdot\Delta f_i$
is precisely the fraction of total wave energy carried by bin $i$. Any
per-frequency log-space error is scaled by exactly how much of the
spectrum's energy that bin represents: a tail bin (small $P(f_i)$)
contributes almost nothing to $D_{KL}$ even if its log-space error is large,
while the same error at the spectral peak (large $P(f_i)$) is weighted
heavily. This is not an arbitrary or hand-tuned weighting scheme — it falls
directly out of treating the spectrum as a probability distribution, and it
reproduces exactly the priority that $H_s$ and $T_{m02}$ already have
physically, since both are dominated by the low-order moments of the
energy-dense region around the peak.

The resulting composite loss:
\begin{equation}
\mathcal{L} = \text{RMSE}\big(E_{\text{pred}}, E_{\text{true}}\big)
\;+\; \lambda_W \cdot W_1\big(\hat P_{\text{pred}}, \hat P_{\text{true}}\big)
\;+\; \lambda_{KL} \cdot D_{KL}\big(\hat P_{\text{true}}\,\|\,\hat P_{\text{pred}}\big)
\end{equation}
gives each term ownership of a distinct failure mode: RMSE anchors absolute
per-bin magnitude (needed for $H_s$/$T_{m02}$ integration fidelity); $W_1$
penalizes energy transported across the frequency axis, including between
distant spectral modes, in proportion to distance moved; $D_{KL}$ penalizes
per-bin log-space error in proportion to physical importance, correcting
the one blind spot RMSE and $W_1$ share (Section~\ref{sec:gap}) without
reintroducing either of their own.

\section{A Further Gap: Local Peak Sharpness}
\label{sec:gap2}

Even with $D_{KL}$ added, RMSE, $W_1$, and $D_{KL}$ together still share one
remaining blind spot: none of them supplies a gradient signal specifically
anchored to whether a detected spectral peak \emph{retains its own
amplitude}. $W_1$'s transport cost is proportional to how far energy moves,
not to how much a peak's own height changes — redistributing mass from a
peak bin into its two immediate neighbors is geometrically cheap (mass moves
only one or two bins) even though it can collapse the peak's amplitude by a
large fraction. $D_{KL}$, as noted in \texttt{SpectralKLDivergenceLoss}'s
own documentation, is dominated by whatever probability the prediction
assigns exactly at the bins where the true distribution concentrates its
mass — this makes it \emph{decline to pile on} for a moderately broadened
peak the way an unweighted MSE would, but it supplies no positive, targeted
gradient pressure to restore the peak's amplitude once it has partially
collapsed; it merely penalizes the collapse less harshly than MSE, rather
than correcting it. Neither mechanism, individually or combined, behaves
like a term whose purpose is specifically "does the peak's height survive."

This matters physically because spectral bandwidth and peakedness — closely
tied to $T_{m02}$ and to distinguishing a narrow, long-travelled swell
partition from a broader, locally generated wind-sea partition — depend on
whether the peak itself stays sharp, not merely on whether energy remains
somewhere in its general neighborhood.

\section{Soft-Max Peak Height as a Fourth Penalty Term}
\label{sec:softmax}

\subsection{From Hard Maximum to a Differentiable Relaxation}

Within a window $\mathcal{W}_k$ around a detected peak $k$, the natural
summary of "the peak's height" is $\max_{i\in\mathcal{W}_k} E(f_i)$. As a
training signal this is unusable: its gradient is exactly zero at every bin
except the single arg-max bin, and that bin can flip discontinuously between
optimization steps. Replacing the hard maximum with a temperature-controlled
soft-max recovers a smooth, everywhere-differentiable surrogate:
\begin{equation}
\sigma_k(i;\tau) = \frac{\exp\big(E(f_i)/\tau\big)}
{\sum_{j\in\mathcal{W}_k}\exp\big(E(f_j)/\tau\big)}, \qquad
H_k^{\text{pred}}(\tau) = \sum_{i\in\mathcal{W}_k} E_{\text{pred}}(f_i)\,\sigma_k(i;\tau).
\end{equation}
As $\tau\to\infty$, $\sigma_k\to$ uniform and $H_k^{\text{pred}}\to$ the
window's plain mean; as $\tau\to0$, $\sigma_k\to$ one-hot at the arg-max and
$H_k^{\text{pred}}\to\max$. For any finite $\tau>0$ every bin in the window
receives nonzero, informative gradient, weighted by its own current
soft-max confidence.

\subsection{Windows from the Existing Partition Detector}

$\mathcal{W}_k$ is taken directly from the trough-to-trough partition
boundaries already computed by this project's Portilla et al. (2009)
significant-peak detector (\texttt{find\_significant\_peaks}), exposed for
this purpose by a new \texttt{find\_peak\_windows} helper in
\texttt{utils/spectral\_partitioning.py} — narrow for a narrow swell
partition, wide for a broad wind-sea partition, rather than an arbitrary
fixed bin-radius that would necessarily be wrong for one of the two regimes.

\subsection{Calibrating the Temperature}

A window's width alone is not a safe basis for $\tau_k$, for two reasons.
First, the buoy's frequency grid is log-spaced and non-uniform (bins
0.005~Hz wide below 0.1~Hz, 0.02~Hz wide above 0.365~Hz), so the number of
bins in a partition is not a reliable proxy for its physical bandwidth — the
window width must be measured as $\Delta f_k = f(r_k)-f(\ell_k)$ in Hz, not
as a bin count. Second, $\exp(E_i/\tau)$ is only numerically and
conceptually meaningful when $\tau$ is expressed in $E$'s own units
(spectral density), not in frequency's — a temperature calibrated purely
from bandwidth would make the soft-max's sharpness swing arbitrarily between
a calm-sea sample and a storm sample with an identically \emph{shaped}
partition, since $E$'s absolute scale carries no relation to a window's
width in Hz. Combining both corrections, the temperature used is a fixed
fraction of the peak's own (label-side) height, modulated by that peak's
bandwidth relative to the grid's total span:
\begin{equation}
\tau_k = \max\!\left(\tau_{\min},\; c \cdot H_k^{\text{true}} \cdot
\frac{f(r_k)-f(\ell_k)}{f_N-f_1}\right),
\end{equation}
so a narrow swell partition gets a small, near-hard-max temperature, and a
broad wind-sea partition gets a larger, softer one — matching the physical
fact that a broad wind-sea peak genuinely does not have as sharply defined a
height as a narrow swell peak. $H_k^{\text{true}}$ is drawn only from the
label, exactly as the window detection itself already is, so this is not a
train/inference information leak.

\subsection{Loss Definition}

On the label side a hard maximum is used directly (no gradient is needed
there): $H_k^{\text{true}} = \max_{i\in\mathcal{W}_k} E_{\text{true}}(f_i)$.
The per-peak, then per-sample, loss is
\begin{equation}
\mathcal{L}_{\text{peak},k} = \big(H_k^{\text{pred}}(\tau_k) -
H_k^{\text{true}}\big)^2, \qquad
\mathcal{L}_{\text{peak}} = \frac{1}{K}\sum_{k=1}^K \mathcal{L}_{\text{peak},k},
\end{equation}
averaged over a sample's $K$ detected peaks (samples with $K=0$ are excluded
from the batch average rather than contributing a filled-in zero, the same
convention already used for undefined-quantity samples via
\texttt{M0\_MASK\_THRESHOLD} in \texttt{nn/evaluate.py}). Adding this term
gives the composite loss its fourth and final component:
\begin{equation}
\mathcal{L} = \text{RMSE}\big(E_{\text{pred}}, E_{\text{true}}\big)
\;+\; \lambda_W\, W_1\big(\hat P_{\text{pred}}, \hat P_{\text{true}}\big)
\;+\; \lambda_{KL}\, D_{KL}\big(\hat P_{\text{true}}\,\|\,\hat P_{\text{pred}}\big)
\;+\; \lambda_{\text{peak}}\, \mathcal{L}_{\text{peak}}.
\end{equation}
$\mathcal{L}_{\text{peak}}$ is implemented and unit-tested as
\texttt{SoftPeakHeightLoss} in \texttt{utils/loss.py}, together with
\texttt{find\_peak\_windows} in \texttt{utils/spectral\_partitioning.py};
neither is yet wired into the training loop's hyperparameter search (the
weight $\lambda_{\text{peak}}$ is not yet an Optuna-tuned parameter), pending
the same promotion path already followed by
\texttt{SpectralWassersteinLoss} and \texttt{SpectralKLDivergenceLoss}.

\section{Wasserstein, KL Divergence, and Soft-Max Localization in Analogous Fields}
\label{sec:literature}

All three mechanisms are well established in other domains, for the same
underlying reason: point-wise metrics ignore the geometry or relative
importance of the axis being compared, and a hard extremum operator (max/
arg-max) carries no usable gradient at all — each of the three corrections
addresses one of these structural gaps.

\subsection{Wasserstein Distance}

\begin{itemize}[leftmargin=1.4em]
\item \textbf{Image retrieval.} Rubner, Tomasi, and
Guibas~\cite{rubner2000} introduced the Earth Mover's Distance as a
similarity metric between image color histograms precisely because
bin-wise $L_1$/$L_2$ distances penalize a color shift between adjacent bins
as heavily as an unrelated color change.
\item \textbf{Audio and music.} Logan and
Salomon~\cite{logan2001} compared musical timbre similarity by representing
each track's spectrum as a Gaussian mixture and comparing mixtures via
Earth Mover's Distance, since a direct Euclidean spectral comparison fails
when the relevant spectral content is concentrated in slightly shifted
frequency regions.
\item \textbf{Biomedical signals / EEG.} Cazelles, Robert, and
Tobar~\cite{cazelles2021} define the ``Wasserstein--Fourier distance''
directly between power spectral densities of stationary time series —
essentially the same problem addressed here, treating the normalized PSD as
a probability distribution over frequency.
\item \textbf{Signal processing, broadly.} Kolouri et
al.~\cite{kolouri2017} survey optimal-transport/Wasserstein methods across
signal processing and machine learning, including spectral and waveform
comparison.
\item \textbf{Deep generative models.} Arjovsky, Chintala, and
Bottou~\cite{arjovsky2017}, in proposing the Wasserstein GAN, show formally
that the Wasserstein distance stays continuous and gradient-informative even
when two distributions have non-overlapping support — directly relevant to
narrow, poorly aligned spectral peaks. Efficient computation is discussed by
Cuturi~\cite{cuturi2013} (Sinkhorn distances) and Peyré and
Cuturi~\cite{peyre2019} (a general computational optimal transport
reference), with a practical implementation available in
POT~\cite{flamary2021}.
\end{itemize}

\subsection{KL Divergence and Weighted Spectral Divergences}

\begin{itemize}[leftmargin=1.4em]
\item \textbf{Information theory foundations.} The divergence itself is due
to Kullback and Leibler~\cite{kullback1951}; its role as the natural
divergence underlying cross-entropy, and its relationship $D_{KL}=H(P,Q)-H(P)$,
is standard material in information theory~\cite{cover2006} and is the same
mechanism underlying cross-entropy loss in classifier training in
general~\cite{goodfellow2016}.
\item \textbf{Audio spectrogram factorization.} Lee and
Seung~\cite{lee2001} introduce a generalized KL-divergence cost function for
non-negative matrix factorization, alongside the Euclidean one, precisely
because it better respects the multiplicative, magnitude-sensitive structure
of non-negative spectral data than a squared-error cost — the same
motivation as Section~\ref{sec:gap}. Févotte, Bertin, and
Durrieu~\cite{fevotte2009} extend this line of work with the closely related
Itakura-Saito divergence for music spectrogram factorization, again choosing
a divergence family that weights spectral errors by magnitude/importance
rather than treating every bin uniformly, for the same physical reason
argued here: a fixed absolute error matters far more where the signal
carries real energy than where it does not.
\end{itemize}

\subsection{Differentiable Soft-Max Localization}

\begin{itemize}[leftmargin=1.4em]
\item \textbf{Human pose estimation.} Sun et
al.~\cite{sun2018} replace the standard, non-differentiable arg-max
decoding of a heatmap-based keypoint detector with an "integral regression"
that computes the joint coordinate as a soft-max-weighted expectation over
the heatmap, precisely to obtain an end-to-end differentiable localization
signal — the same replacement of max/arg-max by a soft-max-weighted average
used here for spectral peak height.
\item \textbf{Discrete relaxation, more generally.} Jang, Gu, and
Poole~\cite{jang2017} formalize the general principle that a temperature-
controlled soft-max is a continuous relaxation of a discrete arg-max
decision, recovering the hard decision as temperature $\to 0$ and a smooth,
fully differentiable surrogate for any finite temperature — exactly the two
limits (mean / hard-max) that $\tau_k$ interpolates between here.
\item \textbf{Attention as soft-max pooling.} Vaswani et
al.~\cite{vaswani2017} popularize soft-max-weighted pooling as the core
mechanism of attention; this project's own encoder already uses the same
mechanism to collapse a spectrum's frequency bins into one embedding
(\texttt{\_FreqAttentionPool}, \texttt{nn/freq\_embedding.py}) — the
soft-max peak-height loss applies the identical mechanism as a loss term
rather than as a model layer.
\end{itemize}

The recurring pattern across these domains — image histograms, audio
spectra, EEG power spectral densities, non-negative spectrogram
factorization, heatmap-based keypoint localization — is the same one
motivating this project's loss design: when the compared object is
fundamentally a \emph{distribution of mass over a structured, physically
meaningful axis} rather than a vector of independent coordinates, the
geometry of that axis (Wasserstein), the relative importance of each
location on it (KL divergence), and a smooth relaxation of any hard
extremum operator on it (soft-max) each carry information that a plain
point-wise metric structurally cannot represent.

\section{Conclusion}

Bin-wise RMSE is necessary but not sufficient for evaluating wave spectra:
it anchors physical magnitude but has no notion of frequency-axis geometry
or of how much any given bin actually matters, producing a known bias
toward smoothed/blurred predictions and blindness to both modality errors
and to systematically undervaluing peak accuracy relative to the spectral
tail. The Wasserstein distance corrects the first gap by penalizing energy
transport in proportion to distance; the KL divergence corrects the second
by weighting each bin's error in proportion to the physical energy it
represents — and, as Section~\ref{sec:kl} shows, is mathematically nothing
more than a $P$-weighted average of the same log-space errors RMSE already
computes, unweighted. A further gap remains even so (Section~\ref{sec:gap2}):
neither mechanism supplies targeted gradient toward restoring a peak's own
amplitude once it has locally collapsed into its neighbors, since $W_1$'s
cost tracks distance moved (small for a local collapse) and $D_{KL}$ only
declines to penalize it as harshly as MSE would, rather than correcting it.
A differentiable soft-max peak-height term (Section~\ref{sec:softmax})
closes this gap directly, replacing the zero-gradient hard maximum with a
temperature-controlled, everywhere-smooth relaxation whose temperature is
itself calibrated from the peak's own physical bandwidth and energy scale.
All three corrections are supported by extensive precedent in analogous
domains facing the same structural problems. The resulting four-term loss
\begin{equation*}
\mathcal{L} = \text{RMSE} + \lambda_W\, W_1 + \lambda_{KL}\, D_{KL}
+ \lambda_{\text{peak}}\, \mathcal{L}_{\text{peak}}
\end{equation*}
gives each term ownership of a distinct, complementary failure mode rather
than asking any single metric to cover all of them at once.
\texttt{SpectralWassersteinLoss}, \texttt{SpectralKLDivergenceLoss}, and
\texttt{SoftPeakHeightLoss} are all implemented and unit-tested in this
project. \texttt{SpectralWassersteinLoss}'s weight is already an
Optuna-tuned hyperparameter; \texttt{SpectralKLDivergenceLoss}'s weight is
wired into the training loop as an opt-in term, pending a manual sweep
before the same promotion; \texttt{SoftPeakHeightLoss} and its supporting
\texttt{find\_peak\_windows} helper are implemented and tested in isolation
but not yet wired into the training loop at all — that integration (
pre-computing per-sample peak windows once, alongside \texttt{freq\_means}/
\texttt{shape\_means}, and threading a $\lambda_{\text{peak}}$ weight
through \texttt{train\_one\_epoch}) is the next step before this term can
be trained against.

\begin{thebibliography}{9}

\bibitem{rubner2000}
Rubner, Y., Tomasi, C., \& Guibas, L. J. (2000).
\textit{The Earth Mover's Distance as a Metric for Image Retrieval}.
International Journal of Computer Vision, 40(2), 99--121.

\bibitem{logan2001}
Logan, B., \& Salomon, A. (2001).
\textit{A Music Similarity Function Based on Signal Analysis}.
IEEE International Conference on Multimedia and Expo (ICME).

\bibitem{cazelles2021}
Cazelles, E., Robert, A., \& Tobar, F. (2021).
\textit{The Wasserstein--Fourier Distance for Stationary Time Series}.
IEEE Transactions on Signal Processing, 69, 709--721.

\bibitem{kolouri2017}
Kolouri, S., Park, S. R., Thorpe, M., Slep\v{c}ev, D., \& Rohde, G. K. (2017).
\textit{Optimal Mass Transport: Signal Processing and Machine-Learning Applications}.
IEEE Signal Processing Magazine, 34(4), 43--59.

\bibitem{arjovsky2017}
Arjovsky, M., Chintala, S., \& Bottou, L. (2017).
\textit{Wasserstein GAN}.
Proceedings of the 34th International Conference on Machine Learning (ICML).

\bibitem{cuturi2013}
Cuturi, M. (2013).
\textit{Sinkhorn Distances: Lightspeed Computation of Optimal Transport}.
Advances in Neural Information Processing Systems (NeurIPS) 26.

\bibitem{peyre2019}
Peyré, G., \& Cuturi, M. (2019).
\textit{Computational Optimal Transport: With Applications to Data Science}.
Foundations and Trends in Machine Learning, 11(5--6), 355--607.

\bibitem{flamary2021}
Flamary, R., Courty, N., Gramfort, A., et al. (2021).
\textit{POT: Python Optimal Transport}.
Journal of Machine Learning Research, 22(78), 1--8.

\bibitem{kullback1951}
Kullback, S., \& Leibler, R. A. (1951).
\textit{On Information and Sufficiency}.
Annals of Mathematical Statistics, 22(1), 79--86.

\bibitem{cover2006}
Cover, T. M., \& Thomas, J. A. (2006).
\textit{Elements of Information Theory} (2nd ed.).
Wiley.

\bibitem{goodfellow2016}
Goodfellow, I., Bengio, Y., \& Courville, A. (2016).
\textit{Deep Learning}.
MIT Press.

\bibitem{lee2001}
Lee, D. D., \& Seung, H. S. (2001).
\textit{Algorithms for Non-negative Matrix Factorization}.
Advances in Neural Information Processing Systems (NeurIPS) 13.

\bibitem{fevotte2009}
F\'{e}votte, C., Bertin, N., \& Durrieu, J.-L. (2009).
\textit{Nonnegative Matrix Factorization with the Itakura-Saito Divergence: With Application to Music Analysis}.
Neural Computation, 21(3), 793--830.

\bibitem{sun2018}
Sun, X., Xiao, B., Wei, F., Liang, S., \& Wei, Y. (2018).
\textit{Integral Human Pose Regression}.
European Conference on Computer Vision (ECCV).

\bibitem{jang2017}
Jang, E., Gu, S., \& Poole, B. (2017).
\textit{Categorical Reparameterization with Gumbel-Softmax}.
International Conference on Learning Representations (ICLR).

\bibitem{vaswani2017}
Vaswani, A., Shazeer, N., Parmar, N., Uszkoreit, J., Jones, L., Gomez, A. N.,
Kaiser, {\L}., \& Polosukhin, I. (2017).
\textit{Attention Is All You Need}.
Advances in Neural Information Processing Systems (NeurIPS) 30.

\end{thebibliography}

\end{document}
